\section{Đánh giá so sánh}
\label{sec:evaluation}

\subsection{Thiết lập và phương pháp}
\label{sec:experimental-setup}

Nhóm benchmark ba database trong cùng một pipeline để giảm sai lệch do môi trường. Kết quả dưới đây là \textbf{snapshot} trên corpus synthetic 10K chunks và 200 queries; vì vậy nó phù hợp để so sánh trong điều kiện kiểm soát, không phải ranking chung cho mọi workload production.

\begin{table}[H]
\centering
\caption{Cấu hình thực nghiệm và fairness protocol.}
\label{tab:experimental-setup}
\begin{tabular}{lp{9.6cm}}
\toprule
\textbf{Thành phần} & \textbf{Chi tiết} \\
\midrule
Hardware & AMD Ryzen 7 5800H, 16 logical CPUs, 14\,GiB RAM, NVMe SSD 512\,GB \\
Runtime & Arch Linux, Docker Compose v2, Ollama local \\
Embedding / LLM & \texttt{nomic-embed-text} 768 chiều; \texttt{qwen2.5:1.5b} cho RAG chat \\
Index params & HNSW $M=16$, $ef\_construction=128$, $ef\_search=64$, Cosine distance \\
Dataset & 10.000 synthetic chunks, seed = 42; 200 golden queries có tag ground-truth \texttt{[CID:xxxx]} \\
Metrics & Recall@1/5/10, MRR, average latency, recall-latency tradeoff, RAM footprint, DX score \\
\bottomrule
\end{tabular}
\end{table}

Recall@K đo tỷ lệ ground-truth xuất hiện trong top K kết quả, còn MRR đo vị trí trung bình của kết quả đúng:
\[
\text{MRR} = \frac{1}{|Q|}\sum_{i=1}^{|Q|}\frac{1}{\text{rank}_i}.
\]
Toàn bộ tham số HNSW được đọc từ cấu hình chung thay vì hardcode riêng cho từng wrapper, giúp so sánh nhất quán hơn.

\subsection{Kết quả snapshot}

\begin{table}[H]
\centering
\caption{Recall@K, MRR và latency trung bình (corpus = 10K, queries = 200).}
\label{tab:accuracy}
\resizebox{\textwidth}{!}{%
\begin{tabular}{lcccccc}
\toprule
\textbf{Engine} & \textbf{Recall@1} & \textbf{Recall@5} & \textbf{Recall@10} & \textbf{MRR} & \textbf{Avg Latency (ms)} & \textbf{Errors} \\
\midrule
Milvus  & 18.0 & 34.0 & 44.0 & 0.2492 & 4.14 & 0 \\
Qdrant  & 3.0  & 7.5  & 9.5  & 0.0467 & 4.83 & 0 \\
Weaviate & 3.0  & 8.0  & 9.5  & 0.0472 & 14.47 & 0 \\
\bottomrule
\end{tabular}}
\end{table}

Trong snapshot này, Milvus đạt Recall@K và MRR cao nhất, đồng thời giữ average latency thấp. Qdrant và Weaviate có Recall@10 tương đương nhau, nhưng Qdrant nhanh hơn trong dense retrieval. Tuy nhiên, cách đọc đúng không phải là ``Milvus luôn tốt nhất'': benchmark này dùng synthetic corpus, single-node deployment và một bộ tham số index cố định.

Recall tuyệt đối thấp chủ yếu đến từ \textit{vector collapse}. Corpus synthetic được sinh từ số lượng template ít; nhiều chunk gần như đồng nghĩa nên embedding model gom chúng vào cùng vùng không gian. Khi đó database có thể tìm đúng cụm ngữ nghĩa nhưng chưa chắc trả đúng ID ground-truth cụ thể. Đây là giới hạn của dataset benchmark, không phải bằng chứng rằng các database có chất lượng retrieval thấp trong mọi corpus thật.

\subsection{Recall-latency tradeoff và filter}

\begin{table}[H]
\centering
\caption{Tradeoff sweep: Recall (\%) / Avg latency (ms) theo top\_k.}
\label{tab:tradeoff}
\resizebox{\textwidth}{!}{%
\begin{tabular}{lcccccc}
\toprule
\textbf{Engine} & \textbf{K=1} & \textbf{K=2} & \textbf{K=5} & \textbf{K=10} & \textbf{K=20} & \textbf{K=50} \\
\midrule
Milvus  & 18.0 / 3.65 & 24.0 / 3.68 & 34.0 / 3.67 & 44.0 / 3.88 & 62.0 / 3.85 & 80.0 / 4.17 \\
Qdrant  & 3.0 / 4.57 & 3.5 / 4.71 & 7.5 / 4.80 & 9.5 / 4.87 & 17.5 / 5.10 & 27.0 / 5.82 \\
Weaviate & 3.0 / 16.19 & 3.5 / 15.77 & 8.0 / 15.39 & 9.5 / 15.49 & 15.5 / 14.96 & 24.5 / 15.66 \\
\bottomrule
\end{tabular}}
\end{table}

Milvus giữ đường tradeoff tốt nhất trên snapshot này: Recall tăng đến 80\% ở K=50 trong khi latency vẫn quanh 4\,ms. Qdrant có latency tăng chậm khi K tăng, cho thấy chi phí biên để lấy thêm kết quả khá nhỏ. Weaviate có latency ổn định quanh 15--16\,ms; kết quả này phản ánh dense retrieval path trong benchmark, chưa thể hiện đầy đủ lợi thế native hybrid search của Weaviate.

Qdrant được đo riêng với filter metadata sau warmup: dense-only baseline đạt 3.97\,ms; equality filter đạt 5.29\,ms; range filter đạt 5.20\,ms; combined filter đạt 5.73\,ms. Overhead +1.2--1.8\,ms cho thấy trong thí nghiệm này payload filter tạo chi phí nhỏ và vẫn phù hợp với RAG cần tenant/ACL filtering.

\subsection{Tài nguyên và Developer Experience}

\begin{table}[H]
\centering
\caption{Tài nguyên và DX trong codebase benchmark của nhóm.}
\label{tab:dx-resource}
\resizebox{\textwidth}{!}{%
\begin{tabular}{lccc}
\toprule
\textbf{Metric} & \textbf{Qdrant} & \textbf{Weaviate} & \textbf{Milvus} \\
\midrule
RAM sau benchmark & $\sim$178\,MB & $\sim$200--300\,MB & $\sim$398\,MB \\
Containers cần thiết & 1 & 1 & 3+ \\
SLOC wrapper & \textbf{232} & 306 & 379 \\
Cyclomatic complexity & \textbf{24} & 59 & 36 \\
Third-party imports & \textbf{2} & 4 & 3 \\
Composite DX score (thấp hơn tốt hơn) & \textbf{190.0} & 307.0 & 290.5 \\
\bottomrule
\end{tabular}}
\end{table}

Qdrant có footprint và DX score tốt nhất trong codebase của nhóm, phù hợp với prototype hoặc team nhỏ. Milvus dùng nhiều tài nguyên và dependency hơn, nhưng đó là chi phí đi kèm kiến trúc phân tán. Weaviate nằm giữa hai cực: triển khai vẫn gọn, nhưng wrapper phức tạp hơn vì schema, module và hybrid query surface.

\subsection{Hạn chế}

Thí nghiệm có bốn giới hạn chính. Thứ nhất, synthetic corpus gây vector collapse nên recall tuyệt đối không đại diện cho tài liệu thật. Thứ hai, benchmark chạy single-node nên chưa đo được lợi thế phân tán của Milvus. Thứ ba, Docker resource limits và warmup/cache có thể ảnh hưởng latency tail. Thứ tư, Weaviate được đánh giá chủ yếu bằng dense retrieval snapshot, trong khi điểm mạnh thiết kế của nó là hybrid retrieval. Vì vậy, kết quả nên được đọc như bằng chứng cho một workload cụ thể, không phải kết luận vĩnh viễn.
